\chapter[Why Ritual Works]{Why Ritual Works Even When Outsiders Do Not Understand}
\section{A Birthday Candle}
Begin with a small object: a birthday candle.

A thin little thing on a cake, burning at most two minutes. Yet those minutes contain strict procedure: lights out, people gathering, the entire prescribed song, no substitute. The birthday person closes their eyes, wishes silently, since speaking makes it ineffective, then blows every candle out at once. Miss one and someone gasps, as though a year's luck escaped.

Honestly answer: who receives the wish? Is there a heavenly address? Why does speaking spoil it---fear of light, or a listening rules committee? Who demanded one breath?

None survives questioning, yet none feels freely cancellable. Omit candles and the birthday seems not quite celebrated; speak the wish and a brief silence requires smoothing over. Perhaps nobody believes in candle magic, yet we still light, blow, and close our eyes.

This is \textbf{ritual}: fixed procedures, repeated acts, particular times and places, whose participants may not explain why but feel deviation is wrong. Birthdays, weddings, funerals, school openings, flags, reunion dinners, concert phone lights: important moments rarely escape it. Its puzzle is producing nothing while working. Blowing a candle brings a wish no nearer fulfilment, yet without it something seems not to have happened. What kind of efficacy is this? Let us examine it.

\section{Three Days in a Mourning Hall}
Come into a scene.

Four in the morning, early winter, an old street in a southern Chinese county town. In a two-storey house, a coffin occupies the central ground-floor room. A square table holds portrait, incense burner, two white candles flickering across the old man's face. Sweet incense, burnt paper money, kitchen ginger soup mingle.

He died the night before last, aged eighty-three. A dozen people remain. His eldest son sits in mourning cloth beside the coffin after over forty sleepless hours. Aunts take turns burning paper, ashes curling like black butterflies. People play cards quietly in a corner; long vigils require company to sit through the night with the dead, everyone locally understands. Neighbours' daughters-in-law volunteer in the kitchen, washing vegetables, cutting meat, setting steamers for dawn's mourners.

Fifteen-year-old granddaughter Xiaoman experiences the whole process for the first time and notices oddities.

Grief has a script. Arrivals stand before the portrait, bow or kneel, relatives return acknowledgement. Aunts' lamentations have rhythm and words, rising for particular guests, subsiding after consolation. Xiaoman first dislikes apparent performance, until her third aunt loses the words mid-lament and breaks into uncontrolled, rhythmless sobs. Script and tears do not conflict; script channels tears like a riverbed.

Everyone has a prescribed place: mourning fabric, procession rank, responsibility for breaking a clay basin. Nobody consults a manual. Knowledge lives in bodies, not books.

The three days are tiring, expensive, troublesome. Offering cigarettes to guests, her uncle hoarsely says the old man economised all his life; the final occasion must not be economised upon.

At dawn the procession leaves, suona horns waking the street. White shoes pressing frost, Xiaoman feels a question grown clearer for three days catch in her throat:

Grandpa cannot feel it. For whom are candles, horns, basin, continuous banquets, three days and nights arranged?

Remember the mourning hall. Most of this chapter answers her.

\section{Why All This Trouble?}
State the conflict without rushing.

If the rite serves the dead, it fails: no banquets seen or horns heard. Believers and unbelievers must acknowledge no verifiable receipt from an afterlife. If it serves the living, why name Grandpa the recipient? Tired, hungry, grieving people could rest, eat, embrace. Why the costly detour through incompletely understood rules for three days?

Go deeper. Nothing seems useful towards a tangible goal: candles illuminate little, broken basins solve nothing, mourning cloth barely warms. Yet fallen candles, unbroken basin, disordered sequence make hearts jolt. Why such concern for correctness in activities producing no practical result?

Harder still, when Xiaoman asks why break the basin, her uncle finally says the rules require it, unable to name their maker. Ritual's oddity is \textbf{authority precisely from unexplained origins}. A purposeful sequence is technique, like repairing a bicycle tyre; an unexplained always-and-must sequence becomes ritual. Vague reasons strengthen authority, almost uniquely among human activities.

Classify the hall's details first. Amid variation, almost every ritual uses four elements:

\textbf{Repetition.} Annual birthday songs, repeated funeral procedures, identical prayer words join this occasion to previous ones and generations. Aunts inherit lamentation from their aunts; shared songs place people in a line larger than themselves.

\textbf{Synchrony.} Simultaneous bows, silence, singing, candle-blowing bypass language to tell bodies we are a whole, not strangers coincidentally sharing a room.

\textbf{Bodily action.} Kneeling, prostrating, bowing, joined palms, raised hands, circling, stepping. Thoughts can evade; bodies less easily. An unusual act, kneeling on concrete, acknowledges before thought that this is no ordinary moment.

\textbf{Special space and time.} A living room becomes a mourning hall, temples, churches, auditoriums, arenas, stages become enclaves cut from ordinary space. New Year's vigil and countdown's final second differ from ordinary nights and seconds. Mark a space-time exceptional and its acts and words follow.

Remember these four. County funerals, Mecca, Olympic arenas, phone screens repeatedly use them. What are they repairing?

\section{Those Who Mock and Those Who Kneel}
Two loud positions persist. Give both full strength, then look elsewhere.

\textbf{First: Ritual is superstition, waste, and chains.}

Defend this side, unfairly dismissed as unfeeling or forgetting ancestors.

Four reasons. Cost: funerals consume months of income, weddings begin marriages in debt; medical care, education, better living seem more practical. Time: three-day vigils or seven-day services are costly in modern life, with compulsory pressure undiminished by hardship. Admission to ignorance: believing unbroken basins unsettle the dead or spoken wishes fail hands life's steering to unseen forces; inauspicious surgical timing has delayed real needs for centuries. Deepest, power's friendship: who demands how many prostrations, receives standing crowds, trains bodily obedience until reasons vanish? Ambitious people cannot ignore technology implanting the sense that refusal is wrong.

Take the case seriously. Its fourth point will recur.

\textbf{Second: Ritual is not excess movement but a necessary container.}

Many situations exceed speech: bereavement makes words light; adulthood reshuffles families beyond a casual announcement; marriage reseats hundreds of relationships. Inner landscapes shift tectonically while ordinary language suits level ground. Ritual supplies a crossing, converting unsayable feelings into doable acts: bow, basin, ring exchange, shared song. You need no perfect words; finish the acts together and they speak.

A deeper defence: ritual does not merely complete the event; it is part of the event. The funeral is not aftermath alone. The manner of farewell shapes the life's ending in collective memory and survivors' continuing life around absence.

With both positions heard, widen the view to humanity's many crossings. Their variety itself speaks.

In early-November Mexico, Day of the Dead altars hold photographs, marigolds, sugar skulls, favourite foods. Families bring food and music to graves through the night. Outsiders wonder at laughter. Locals explain not death disguised as celebration but relatives visiting from another world, deserving hospitality. Mourning and reunion share one rite.

At Varanasi's Ganges steps, cremation continues year-round. Families carry colourful cloth-wrapped bodies through lanes, complete public last rites, release ashes into water. Believers see extraordinary blessing in dying here; observers also see the water source supporting millions' drinking, bathing, cultivation. One fire meets two visions.

Consider fasting. Ramadan brings worldwide dawn-to-sunset abstention from food and water, then family and community meals. A growing teenager's first complete fast may be remembered less for hunger than twilight silence awaiting the call to prayer, hundreds of millions waiting for the same sunset. Jewish Yom Kippur includes roughly a day and night without food, confronting yearly debts to others. Fasting acts by stopping, removing food, amusement, ordinary life so the extraordinary appears.

Pilgrimage brings around two million to Mecca in the pilgrimage month, men in two seamless white cloths, king and labourer, rich and poor levelled. Europe's routes to Santiago have carried pilgrims over a millennium; unbelievers still walk hundreds of kilometres, saying the journey clarified things.

Seasonal rites surround you too. Reunion-dinner dishes may change, the family's shared table that night must not. Qingming offerings vary north and south, standing briefly where someone rests does not. Japanese Obon welcome fires, dances, farewell fires annually affirm that family has not dissolved; some members merely moved farther away.

Beware an easy inference that ritual belongs chiefly to scientifically uninformed people who worship before sailing because they lack weather knowledge. Facts trip it up.

Over a century ago, Malinowski observed Trobriand lagoon fishing with little magic versus open-sea journeys with elaborate rites. Crucially, sailors and boatbuilders possessed precise wind, current, fish knowledge. \textbf{Controllable work attracted less ritual; uncontrollable hazards attracted more}. Ritual responds to uncertainty, not replaces knowledge. Before mocking, inspect stock traders' opening routines, players touching posts three times before penalties, candidates' lucky pens, all among scientifically educated people.

Another reverse error claims belief is required. Have you seen an atheist wedding fall silent as a father gives his daughter's hand, or an easygoing boy tear up at graduation? No divine rings or magical caps are needed. Efficacy does not live wholly in belief. Where, then?

\section[Thinking Bridge: Four Layers of Ritual]{Thinking Bridge: Divide Ritual into Four Layers}
Can any rite, foreign or familiar but inarticulate, be analysed without relying on intuition? Four portable questions separate layers.

\textbf{First: What did they do formally?}

Record like a camera, without explanation: time, people, clothing, position, actions, repetitions. Xiaoman's hall: three days, close relatives in mourning, portrait incense and candles, guests bowing and relatives responding, age-ordered procession, eldest son breaking basin. This factual foundation is shared by observers.

\textbf{Second: What functions did it perform?}

What observable changes followed? Long-separated relatives reunite, divide work, share tasks. Grief enters a schedule of paper-burning, meals, vigil shifts, supporting people through early days. The old man's life is recounted and place in communal memory confirmed. No spirit claims are needed for visible comfort, reunion, transition, identity.

\textbf{Third: What did participants understand it to mean?}

The uncle sees sending the dead onward, aunts money reaching another world, Xiaoman's mother perhaps not literal belief but the children's inability to live with an inadequate farewell. Record honestly, neither decorating nor mocking. Meaning is what it means to participants, but belief is not automatically fact.

\textbf{Fourth: What does it claim is true about the world?}

Afterlife, watching ancestors, listening gods, received wishes: claims behind rites may be true or false, disputable through evidence under the same scrutiny as Earth's orbit or vaccines preventing disease.

Separate layers and quarrels clarify. Rain-making is ignorant versus meaningful can last a century because speakers differ in layer. Dance, offerings, sequence are factual. Shared action during drought distributes helplessness and sustains community observably. Belief in the Dragon King's hearing is recorded. A causal claim that dancing produces rain is false, meteorologically testable. Conclusions coexist: \textbf{a false fourth layer does not invalidate the second}. Bad meteorology may be real social first aid. Conversely, comfort does not prove afterlife; a consoling funeral works either way.

Name the tool yourself; its instruction is \textbf{ask which layer before arguing}. Anthems, astrology, lucky carp, forwarded Confucius blessings also divide. Ancient material shows this was mastered long before today.

\section{Xunzi Watches a Rain Rite}
Read Xunzi, the famously sober late Warring States thinker, directly addressing rain-making in ``Discourse on Heaven'':

\begin{quote}
It rains after a rain sacrifice. Why? No special reason: just as it rains without sacrifice.
\end{quote}
The Chinese \textit{yu} denotes rain sacrifice. Xunzi continues: drums rescuing an eclipsed sun, drought rites, divination in uncertainty are not meant, by informed performers, literally to change heaven. They give human affairs cultural form. Then his sharpest line:

\begin{quote}
The gentleman regards it as cultural form; common people regard it as divine. As cultural form it is auspicious; as divine it is disastrous.
\end{quote}
Broadly, informed people see human culture, others divine action; the first is good, the second problematic.

Use four layers. He takes formal practice seriously, glimpses its calming function, distinguishes elite and popular meanings, and rejects supernatural causation: dancing does not bring rain, eclipses are not heavenly dogs eating suns. Twenty-three centuries ago, without satellites, he separated ritual usefulness from its claims' truth.

Another old text, ``Eight Rows'' in the \textit{Analects}:

\begin{quote}
Sacrifice as though present; sacrifice to spirits as though they were present. The Master said: ``If I do not participate, it is as though I had not sacrificed.''
\end{quote}
Ancestors and spirits are treated as before oneself; nonparticipation amounts to no rite.

The subtle word is \textit{ru}, as though. Confucius does not simply assert presence, leaving doors for faith and doubt: literal ancestors or inward presence. The latter clause requires \textbf{bodily attendance}. Substitution and remote gestures do not count. Like the uncle personally keeping vigil and eldest son personally breaking the basin, ritual cannot be delivered or outsourced.

A third text, ``Questions about Three-Year Mourning'' in the \textit{Book of Rites}, explains parental mourning through four Chinese characters:

\begin{quote}
Measure feeling and establish its form.
\end{quote}
Create forms proportionate to inward emotion. Their thickness corresponds to feeling's weight; not arbitrary, but measured against human hearts.

Together, a mature ancient layered analysis emerges. Xunzi separates effects and supernatural claims; Confucius stresses sincere bodily participation; the \textit{Rites} makes form feeling's vessel, not an external chain. Neither abolish fraud nor obey divine access: ritual is a human-made tool, whose use requires understanding.

\section{One Candle, Three Explanations}
With tools and examples, compare efficacy explanations. Separate observed kneeling, singing, fasting, circling; causal explanations; participant meanings such as sending Grandpa onward; and later judgements about retention. We compare the causal layer.

\textbf{First: Ritual soothes individuals (psychology).}

Its client is the person facing lost control, upheaval, inexpressibility. Fishers cannot govern waves but can govern departure procedure, retaining agency. Psychology likewise finds stressed people spontaneously increasing small repetitive rites that stabilise feeling. Bereavement temporarily disables action; schedules decide times and tasks, compliance itself offering support. Wordless emotion finds movement: a prostration weighs more than condolences. Concrete and testable, this explains personal effects. Its limit: why almost universal collective, synchronised digestion of grief rather than solitary consolation?

\textbf{Second: Ritual produces the group (sociology).}

Durkheim's study of Australian totemic rites in \textit{The Elementary Forms of Religious Life} contrasts dispersed daily labour with gathered eating, dancing, crying, laughing. Emotions amplify like currents into collective effervescence, an intense whole. The worshipped god partly embodies the group, society bowing each member before us. Secular graduation, parades, New Year songs similarly join coincidental individuals into a bounded, emotional community with a past. This explains synchrony, gathering, repetition. Its limit: us creates them. Fasting identifies nonfasters; uniform-shirted supporters sing together and boo opponents. Unity and boundary are two faces of one stamp.

\textbf{Third: Ritual serves order and power (politics).}

Coolly ask who benefits most: often whoever stands onstage. Coronation and heaven-worship demonstrate heavenly authority; standing for judges reaffirms legal solemnity; Monday flag speeches distinguish speakers from listeners. Bodies learn power's map without words. Who demands lavish final farewells? Ceremonial businesses, face-conscious culture, clan opinion each play roles. This penetrates interests behind tradition, but cannot wholly explain sincere involvement. The uncle's red eyes and aunt's broken lament are not pretence. Reducing all ritual to manipulation resembles reducing all gifts to exchange: penetrating, incomplete.

Psychology sees hearts, sociology groups, politics their hierarchy. A funeral simultaneously comforts particular grief, joins family, ranks elders and juniors. No bland compromise: complete explanations often achieve completeness by discarding half the rite.

\section[Further Challenge: At the Threshold]{Further Challenge: People on the Threshold}
Now a theory gathering scattered phenomena. Early twentieth-century folklorist Arnold van Gennep found in \textit{The Rites of Passage} a shared three-part structure beneath diverse initiations, weddings, funerals: \textbf{separation, liminality, incorporation}.

Initiation begins by removing children from their group: home to camp, shaved hair, changed clothes, old self cut away. Liminality, threshold existence, follows: neither child nor adult, old identity removed, new not yet worn. Trials, secret knowledge, hunger, sleeplessness often accompany it. Completion returns them differently placed, through feast, name, clothes, public adult recognition.

The reach is remarkable. Brides separate from natal households, occupy gaps between families under veils or carried over thresholds without touching ground, then join new identities. Funerals place both dead and relatives in transition, hence first-seven-day observances and mourners temporarily withdrawing from ordinary greetings and feasts until reintegration. Ancient \textit{Rites} adulthood ceremonies, men capped and given courtesy names at twenty, women hair-pinned at fifteen, dramatise incorporation. Three increasingly solemn caps and changed address rewrite a social cell.

Mid-century Victor Turner deepened liminality. Threshold people temporarily leave rank tables, making unusual equality and intimacy, communitas, possible. Rich sons and farm boys share recruit haircuts; rulers and taxi drivers circle Mecca in white; graduation classmates embrace after years barely speaking. Ordinary society is a ranked seating plan, liminality a brief darkness. Lights return and seats resume, but shoulder-to-shoulder feeling remains. Society needs both arrangement and darkness; ritual is the timed switch.

Weigh limits. Derived from passages, the structure fits initiations, weddings, funerals more tightly than festivals: whom does a countdown separate or incorporate? And description does not judge goodness. Humiliating, frightening initiations deserve criticism despite standard structure. Theory illuminates, not absolves.

\section{Glow Sticks and Rallying Assemblies}
Return to daily life. Ritual is not confined to temples, villages, antiquity. Your age is saturated; temples have changed address.

School supplies the complete set. Opening speeches, student pledges, class-block seating; weekly flags, anthem, thousands looking up. Four elements: weekly repetition with unchanged music, synchrony, attentive bodies, exceptional playground time. Hundred-day examination rallies intensify separation through banners, slogans, fists, cries, extracting ordinary students into special candidates on a threshold. Effective? Many feel briefly inspired. Problematic? When enthusiasm becomes compulsory performance, silent students become suspect.

Concerts seem less ritualistic, yet rhythmic glow sticks turn dark stadiums into breathing light. Durkheim would recognise effervescence with singer replacing deity. New Year crowds count down and embrace strangers; physically ordinary midnight becomes a doorway through ritual. Championship streets filled with team colours and strangers' high-fives summon communitas.

Networks mass-produce rather than abolish ritual. Candle emojis form online mourning halls. Game check-ins create tiny devotion through repetition and unbroken streaks. Birthday messages at midnight make a minute late insufficiently caring, inheriting exact timing. The popularity of ritual feeling reveals people professing only science and efficiency deliberately marking ordinary days exceptional.

Power remembers too: corporate morning chants, uniform team-building, immovable pledges. When synchrony trains obedience and nonparticipation means disloyalty, reread Xunzi. Ritual settles or recruits people; the difference lies in \textbf{understanding what happens and having freedom not to be absorbed}. Distinguishing crossing-place from rein is daily homework.

\section[Your Turn: If Ritual Is Only Useful]{Your Turn: What If Ritual Is Merely Useful?}
Return to candle and mourning hall.

Wishes have no recipient, basins no inspector, and Xunzi saw through this long ago. Yet Xiaoman's tears, her uncle's reddened eyes, your tight throat at graduation songs are real.

Answer the final question with reasons, without a standard solution:

\textbf{If psychological and social mechanisms entirely explain a ritual's comfort, unity, and marked transitions, is it genuinely effective or sophisticated self-deception?}

Weigh both sides. Factual claims face firm standards: moving rain dances do not bring rain, and supernatural promises disguised as comfort have deceived and harmed. Yet embraces are merely arms enclosing bodies. Merely can analyse sharply or reveal emotional shortsightedness. Someone knowingly making an undelivered candle wish may be deceived, or carefully marking life special.

Go deeper. If science explains every brain region, hormone, and increase in cohesion, will ritual fail like exposed conjuring, or remain compelling like a sunset whose causes are understood?

What is your answer, and why?
